%! Author = sriadityavedantam
%! Date = 3/30/26

\begin{problem}{
    Prove that:

    (a)
    \[
        (n-a)^2\equiv a^2 \mod n.
    \]

    (b)
    \[
        (2n-a)^2\equiv a^2 \mod {4n}.
    \]
}
    \textbf{(a).}
    By the definition of congruence, we must show that
    \[
        n\mid a^2-(n-a)^2.
    \]
    Compute:
    \begin{align*}
        a^2-(n-a)^2
        &= a^2-(n^2-2an+a^2)\\
        &= 2an-n^2\\
        &= n(2a-n).
    \end{align*}
    Since the right-hand side is a multiple of $n$, we conclude that
    \[
        (n-a)^2\equiv a^2\mod n.
    \]

    \textbf{(b).}
    By the definition of congruence, we must show that
    \[
        4n\mid a^2-(2n-a)^2.
    \]
    Compute:
    \begin{align*}
        a^2-(2n-a)^2
        &= a^2-(4n^2-4an+a^2)\\
        &= 4an-4n^2\\
        &= 4n(a-n).
    \end{align*}
    Since the right-hand side is a multiple of $4n$, we conclude that
    \[
        (2n-a)^2\equiv a^2\mod {4n}.
    \]
\end{problem}

\begin{problem}{
    If $p\geq 5$ and $p$ is prime, prove that $[p]=[1]$ or $[p]=[5]$ in $\mathbb{Z}_6$.
}
    Theorem 2.3 states that
    \[
        [a]=[b]
        \quad\text{if and only if}\quad
        a\equiv b\mod n.
    \]
    Since we are in $\mathbb{Z}_6$, this means
    \[
        [p]=[1]
        \quad\text{if and only if}\quad
        p\equiv 1\mod 6,
    \]
    and
    \[
        [p]=[5]
        \quad\text{if and only if}\quad
        p\equiv 5\mod 6.
    \]    Also,
    \[
        \mathbb{Z}_6=\{[0],[1],[2],[3],[4],[5]\}.
    \]
    We prove that $[p]$ cannot be any of $[0],[2],[3],[4]$.
    Since $p$ is prime and $p\geq 5$, $p$ is odd.
    Therefore $p$ cannot be congruent to $2$ or $4$ modulo $6$, because those are even classes.
    Also, $p$ cannot be congruent to $0$ modulo $6$, since that would mean
    \[
        6\mid p,
    \]
    which is impossible for a prime $p\geq 5$.
    If $[p]=[3]$, then
    \[
        p\equiv 3\mod 6,
    \]
    so
    \[
        p=6k+3=3(2k+1)
    \]
    for some $k\in\mathbb{Z}$.
    Thus $3\mid p$.
    Since $p$ is prime and $p\geq 5$, this is impossible.
    Therefore $[p]$ cannot be $[0],[2],[3]$, or $[4]$ in $\mathbb{Z}_6$.
    Hence the only remaining possibilities are
    \[
        [p]=[1]
        \quad\text{or}\quad
        [p]=[5].
    \]
\end{problem}

\begin{problem}{
    Prove that
    \[
        10^n\equiv (-1)^n\mod {11}
    \]
    for every positive integer $n$.
}
    We prove this using mathematical induction.
    For the base case, let $n=1$.
    Then
    \begin{align*}
        10^1 &\equiv (-1)^1\mod {11}\\
        10 &\equiv -1\mod {11},
    \end{align*}
    which is true since
    \[
        11\mid (-1)-10=-11.
    \]
    Now define
    \[
        S\coloneqq \{x\in\mathbb{Z}^+ : 10^x\equiv(-1)^x\mod {11}\},
    \]
    and we will show that
    \[
        S=\mathbb{Z}^+.
    \]
    Assume $n\in S$.
    Then
    \[
        10^n\equiv (-1)^n\mod {11}.
    \]
    We must show that
    \[
        10^{n+1}\equiv (-1)^{n+1}\mod {11}.
    \]
    Now
    \[
        10^{n+1}=10\cdot 10^n.
    \]
    Also,
    \[
        (-1)^{n+1}=(-1)\cdot (-1)^n.
    \]
    Since
    \[
        10\equiv -1\mod {11}
    \]
    and
    \[
        10^n\equiv (-1)^n\mod {11},
    \]
    Theorem 2.2 implies that
    \[
        10\cdot 10^n\equiv (-1)\cdot (-1)^n\mod {11}.
    \]
    So
    \[
        10^{n+1}\equiv (-1)^{n+1}\mod {11}.
    \]
    Thus
    \[
        n+1\in S.
    \]
    Therefore, by the Principle of Mathematical Induction,
    \[
        S=\mathbb{Z}^+.
    \]
    Hence
    \[
        10^n\equiv (-1)^n\mod {11}
    \]
    for every positive integer $n$.
\end{problem}

\begin{problem}{
    (a) Prove or disprove: If
    \[
        a^2\equiv b^2\mod n,
    \]
    then
    \[
        a\equiv b\mod n
        \quad\text{or}\quad
        a\equiv -b\mod n.
    \]

    (b) Do part (a) when $n$ is prime.
}
    \textbf{(a).}
    This statement is false in general.
    A disproof by example is given by
    \[
        a=6,\qquad b=1,\qquad n=35.
    \]
    Then
    \[
        6^2=36\equiv 1=1^2\mod {35}.
    \]
    So
    \[
        a^2\equiv b^2\mod {35}.
    \]
    However,
    \[
        6\not\equiv 1\mod {35}
    \]
    and
    \[
        6\not\equiv -1\mod {35}.
    \]
    Thus the statement is false.

    \begin{remark}
        I have found some commonality for part (a).
        I noticed that the product of twin primes can equal $n$, and $a$ and $b$ are some sum and difference between the two, generally $a$ being one less than the higher twin and $b=1$.
        From what I have calculated, this seems to be the general idea, but I have not proven it for all twin primes.
        So it is still a conjecture.
        \begin{gather*}
            a=4,\ b=1,\ n=15\\
            a=12,\ b=1,\ n=143\\
            a=18,\ b=1,\ n=323\\
        \end{gather*}
        And so forth.
    \end{remark}

    \textbf{(b).}
    Now assume that $n$ is prime.
    Since
    \[
        a^2\equiv b^2\mod n,
    \]
    we have
    \[
        n\mid b^2-a^2.
    \]
    Factor:
    \[
        b^2-a^2=(b-a)(b+a).
    \]
    So
    \[
        n\mid (b-a)(b+a).
    \]
    Since $n$ is prime, Euclid's Lemma implies that
    \[
        n\mid (b-a)
        \quad\text{or}\quad
        n\mid (b+a).
    \]
    If
    \[
        n\mid (b-a),
    \]
    then
    \[
        a\equiv b\mod n.
    \]
    If
    \[
        n\mid (b+a),
    \]
    then
    \[
        a\equiv -b\mod n.
    \]
    Thus, when $n$ is prime,
    \[
        a^2\equiv b^2\mod n
        \implies
        a\equiv b\mod n
        \quad\text{or}\quad
        a\equiv -b\mod n.
    \]
\end{problem}

\begin{problem}{
    As we have clearly stated in class, and given the defined answer for parts (a) and (b), prove the following.

    (a) If $a$ is a unit in $\zn$, prove that $a$ is not a zero divisor.

    (b) If $a$ is a zero divisor in $\zn$, prove that $a$ is not a unit.
}
    \textbf{(a).}
    Assume, by the definition of a unit, that
    \[
        \gcd(a,n)=1.
    \]
    Suppose there is a $b\in\mathbb{Z}$ with
    \[
        [a][b]=[0].
    \]
    To show that $[a]$ is not a zero divisor, we must prove that this implies
    \[
        [b]=[0].
    \]
    By the previous theorem, $[a]$ is a unit in $\zn$, so there exists $x\in\mathbb{Z}$ such that
    \[
        [x][a]=[1].
    \]
    Then
    \begin{align*}
        [x]([a][b]) &= [x][0]\\
        &= [x\cdot 0]\\
        &= [0].
    \end{align*}
    By associativity,
    \begin{align*}
        ([x][a])[b] &= [0]\\
    [1][b] &= [0]\\
    [b] &= [0].
    \end{align*}
    Hence $[a]$ is not a zero divisor.

    \textbf{(b).}
    Conversely, suppose that
    \[
        \gcd(a,n)=d>1.
    \]
    Then
    \[
        d\mid a
        \qquad\text{and}\qquad
        d\mid n.
    \]
    So
    \[
        a=d\left(\frac{a}{d}\right)
        \qquad\text{and}\qquad
        n=d\left(\frac{n}{d}\right),
    \]
    where
    \[
        \frac{a}{d},\frac{n}{d}\in\mathbb{Z}.
    \]
    Now
    \[
        [a]\left[\frac{n}{d}\right]
        =
        \left[a\cdot \frac{n}{d}\right]
        =
        \left[d\cdot \frac{a}{d}\cdot \frac{n}{d}\right]
        =
        \left[n\cdot \frac{a}{d}\right]
        =
        [0].
    \]
    Since $d>1$, we have
    \[
        \left[\frac{n}{d}\right]\neq [0]
    \]
    in $\zn$.
    Thus $[a]$ is a zero divisor.
    Therefore, if $a$ is a zero divisor in $\zn$, then $a$ is not a unit.
\end{problem}

\begin{problem}{
    Prove that every nonzero element of $\zn$ is either a unit or a zero divisor, but not both.
}
    Let $[a]\neq [0]$ in $\zn$.
    If
    \[
        \gcd(a,n)=1,
    \]
    then by the previous theorem $[a]$ is a unit.
    If
    \[
        \gcd(a,n)\neq 1,
    \]
    then since $[a]\neq [0]$, we have
    \[
        \gcd(a,n)=d>1,
    \]
    and by the previous problem $[a]$ is a zero divisor.
    So every nonzero element of $\zn$ is either a unit or a zero divisor.
    Now we show it cannot be both.
    Suppose $[a]$ is both a unit and a zero divisor.
    Then there exist $[b],[c]\in\zn$ with
    \[
        [b]\neq [0],
        \qquad
        [a][b]=[0],
        \qquad
        [a][c]=[1].
    \]
    Then
    \begin{align*}
        ([a][b])[c] &= [0][c]=[0].
    \end{align*}
    By associativity,
    \begin{align*}
        ([a][c])[b] &= [1][b]=[b].
    \end{align*}
    But the left-hand sides are equal by associativity and commutativity of multiplication in $\zn$, so we get
    \[
        [b]=[0],
    \]
    a contradiction.
    Therefore every nonzero element of $\zn$ is either a unit or a zero divisor, but not both.
\end{problem}

\begin{problem}{
    Find the multiplicative inverse of $[9]$ in $\mathbb{Z}_{41}$.
}
    We want to find $[x]$ such that
    \[
        [9][x]=[1].
    \]
    This is the same as finding integers $x$ and $q$ such that
    \[
        9x=1+41q,
    \]
    or equivalently,
    \[
        9x-41q=1.
    \]
    Since
    \[
        \gcd(41,9)=1,
    \]
    there exist integers $x$ and $y$ such that
    \[
        9x+41y=1.
    \]
    We find them using the Euclidean algorithm:
    \begin{align*}
        41 &= 4\cdot 9 + 5,\\
        9 &= 1\cdot 5 + 4,\\
        5 &= 1\cdot 4 + 1,\\
        4 &= 4\cdot 1 + 0.
    \end{align*}
    Now work backward:
    \begin{align*}
        1 &= 5-1\cdot 4\\
        &= 5-(9-1\cdot 5)\\
        &= 2\cdot 5-9\\
        &= 2(41-4\cdot 9)-9\\
        &= 2\cdot 41-9\cdot 9.
    \end{align*}
    So
    \[
        (-9)9 + 2\cdot 41 = 1.
    \]
    Thus
    \[
        [-9][9]=[1]
    \]
    in $\mathbb{Z}_{41}$.
    Since
    \[
        -9\equiv 32\mod {41},
    \]
    the multiplicative inverse of $[9]$ in $\mathbb{Z}_{41}$ is
    \[
        [32].
    \]
\end{problem}